\section*{Erd\H{o}s Problem \#354: two floor-doubling sequences}
\subsection*{1. Statement and scope}
Given $\alpha,\beta>0$ with $\alpha/\beta$ irrational, must the indexed
multiset
\[
 \{\lfloor2^s\alpha\rfloor:s\ge0\}
 \ \cup\ \{\lfloor2^t\beta\rfloor:t\ge0\}
\]
be complete? Each index can be used once; equal values at different
indices are separate available terms. Completeness means that every
sufficiently large integer is a finite subset sum.
There is also the variant with base $\gamma\in(1,2)$.
We give a full elementary proof of Hegyv\'ari's known
dyadic--nondyadic case and an exact criterion when both parameters
are dyadic. The general irrational-ratio question remains unresolved.

\subsection*{2. Residues of subset sums of one nondyadic sequence}
Let $\beta$ be nondyadic and write $b_j=\lfloor2^j\beta\rfloor$.
Then
\[
 b_{j+1}=2b_j+\varepsilon_j,\qquad\varepsilon_j\in\{0,1\}.
\]
There are infinitely many indices with $\varepsilon_j=1$.
Otherwise $b_{j+1}=2b_j$ eventually, and
$\beta=\lim 2^{-j}b_j$ would be dyadic.
It follows that for every modulus $m\ge1$ and every $J$,
\[
 \gcd(m,b_J,b_{J+1},\ldots)=1.                    \tag{1}
\]
Indeed, any common divisor of the displayed numbers divides every
tail increment $\varepsilon_j$, including an increment equal to $1$.

Let $R_N\subseteq\mathbb Z/m\mathbb Z$ be the residues of subset sums
of $b_J,\ldots,b_N$, including $0$. These finite sets increase with $N$,
so eventually stabilize to a nonempty set $R$. At every subsequent
index $j$, adjoining $b_j$ gives
$R\cup(R+b_j)=R$, hence $R+b_j=R$ by equality of finite cardinalities.
Therefore $R$ is invariant under the subgroup generated by all tail
$b_j$. By (1) this is the entire cyclic group, so $R=\mathbb Z/m\mathbb Z$.
Every residue is thus achieved by a subset sum from a finite prefix
of \emph{any} tail. No assumption that single terms hit every residue
is needed.

\subsection*{3. Complete reconstruction of the dyadic--nondyadic case}
Write $\alpha=m/2^s$ with $m\ge1$ an integer and $s\ge0$.
The $\alpha$-sequence contains the tail
$m,2m,4m,\ldots$, whose subset sums are exactly $m\mathbb Z_{\ge0}$.
By the preceding lemma choose, for each residue $r$ modulo $m$,
a finite $\beta$-subset sum $B_r\equiv r\pmod m$. Set $C=\max_r B_r$.
For every $N\ge C$, choose $r\equiv N\pmod m$ and write
\[
 N=B_r+m q,\qquad q\ge0.
\]
The second summand is a binary subset sum of the $\alpha$-tail.
The two index families are disjoint, so no term is reused.
This proves completeness whenever one parameter is dyadic and the
other is not, including the corresponding irrational-ratio cases.

\subsection*{4. An exact finite criterion when both parameters are dyadic}
Suppose $\alpha=a/2^s$ and $\beta=b/2^t$ with positive integers $a,b$.
Let $H$ be the finite multiset of the terms before the respective
tails $a,2a,\ldots$ and $b,2b,\ldots$, and put $g=\gcd(a,b)$.
Then
\[
 P(\alpha,\beta)=P(H)+a\mathbb Z_{\ge0}+b\mathbb Z_{\ge0}.            \tag{2}
\]
The numerical semigroup on the right contains every sufficiently
large multiple of $g$: dividing by $g$, write $A=a/g,B=b/g$;
for each residue modulo $A$ choose $0\le v<A$ with $Bv$ in that
residue. Every sufficiently large integer in the residue is
$Bv+Au$ with $u\ge0$.
Consequently (2) is complete if and only if the subset sums of
$H$ hit every residue modulo $g$. Necessity follows because the
two tails contribute only multiples of $g$; sufficiency follows
by choosing a fixed head sum in each residue and then using the tails.
The criterion is independent of the chosen dyadic representations,
because (2) always gives exactly the same subset-sum set.

\subsection*{5. Assessment and provenance}
Both arguments use an exact tail supplying all sufficiently large
multiples of a fixed modulus. For two general nondyadic parameters,
the residue lemma remains valid, but it does not supply such a tail.
Surjectivity modulo every modulus is not itself completeness.
The general irrational-ratio case and the general $\gamma\in(1,2)$
variant remain unresolved in this attempt.

The statement and attribution of the special case to Hegyv\'ari [He89]
are at \texttt{https://www.erdosproblems.com/354}, checked 2026-09-25.
The earlier GPT Pro 5.2 special-case argument was checked and
reconstructed; in particular one must not infer $\gcd(R)=1$ merely
from $\gcd(m,R)=1$. The finite-group argument above avoids that
incorrect inference. All inputs to the proof are elementary and
proved here. No novelty or formal verification is claimed.
\subsection*{COMPLETION ESTIMATE}
\textbf{COMPLETION: 60\%}
